\documentclass[12pt]{article}
\usepackage[a4paper,margin=1in]{geometry}
\usepackage[utf8]{inputenc}
\usepackage[T1]{fontenc}
\usepackage{amsmath,amssymb}
\usepackage{textcomp}
\usepackage{parskip}
\begin{document}
\begin{center}
{\Large \textbf{2014 STEP III}}
\end{center}

\subsection*{Page 1}
\begin{flushleft}
\bigskip\textbf{Section A: Pure Mathematics}\par
\bigskip\textbf{Question 1.} Let a, b and c be real numbers such that a + b + c = 0 and let\par
\[(1 + ax)(1 + bx)(1 + cx) = 1 + qx2 + rx3\]
\[for all real x. Show that q = bc + ca + ab and r = abc.\]
\textbf{(i)} Show that the coefficient of xn in the series expansion (in ascending powers of x) of\par
\[ln(1 + qx2 + rx3) is ( 1)n+1S where\]
n\par
-\par
an + bn + cn\par
\[S = , (n ⩾ 1).\]
n\par
n\par
\textbf{(ii)} Find, in terms of q and r, the coefficients of x2, x3 and x5 in the series expansion (in\par
\[ascending powers of x) of ln(1 + qx2 + rx3) and hence show that S S = S .\]
2 3 5\par
\textbf{(iii)} Show that S S = S .\par
2 5 7\par
\textbf{(iv)} Give a proof of, or find a counterexample to, the claim that S S = S .\par
2 7 9\par
\bigskip\textbf{Question 2.} (i) Show, by means of the substitution u = cosh x , that\par
\[sinh x 1 \sqrt\{2\} cosh x 1\]
\[dx = ln - + C .\]
\{ \{\par
\[∫ cosh 2x 2\sqrt\{2\} \{ \sqrt\{2\} cosh x + 1 \{\]
\{ \{\par
\{ \{\par
\{ \{\par
\textbf{(ii)} Use a similar substitution to find an expression for\par
cosh x\par
dx .\par
∫ cosh 2x\par
\textbf{(iii)} Using parts (i) and (ii) above, show that\par
\[1 1 \pi + 2 ln(\sqrt\{2\} + 1)\]
\[du = .\]
\[∫ 1 + u^4 4\sqrt\{2\}\]
\end{flushleft}
\subsection*{Page 2}
\begin{flushleft}
\bigskip\textbf{Question 3.} (i) The line L has equation y = mx + c, where m > 0 and c > 0. Show that, in the case\par
\[mc > a > 0, the shortest distance between L and the parabola y^2 = 4ax is\]
mc a\par
- .\par
\[m\sqrt\{m^2\} + 1\]
What is the shortest distance in the case that mc ⩽ a?\par
\textbf{(ii)} Find the shortest distance between the point (p, 0), where p > 0, and the parabola\par
\[y^2 = 4ax, where a > 0, in the different cases that arise according to the value of p/a.\]
\[[You may wish to use the parametric coordinates (at2, 2at) of points on the parabola.]\]
\[Hence find the shortest distance between the circle (x p)2 + y^2 = b^2, where p > 0\]
-\par
\[and b > 0, and the parabola y^2 = 4ax, where a > 0, in the different cases that arise\]
according to the values of p, a and b.\par
\bigskip\textbf{Question 4.} (i) Let\par
1 1\par
\[I = (y ′ )2 y^2 dx and I = (y ′ + y tan x)2 dx ,\]
∫ - ∫\par
0 ( ) 0\par
\[where y is a given function of x satisfying y = 0 at x = 1. Show that I I = 0 and\]
-\par
\[deduce that I ⩾ 0. Show further that I = 0 only if y = 0 for all x (0 ⩽ x ⩽ 1).\]
\textbf{(ii)} Let\par
\[J = (y ′ )2 a2y2 dx ,\]
∫ -\par
0 ( )\par
where a is a given positive constant and y is a given function of x, not identically zero,\par
\[satisfying y = 0 at x = 1. By considering an integral of the form\]
(y ′ + ay tan bx)2 dx ,\par
∫\par
where b is suitably chosen, show that J ⩾ 0. You should state the range of values of a,\par
\[in the form a < k, for which your proof is valid.\]
\[In the case a = k, find a function y (not everywhere zero) such that J = 0.\]
\end{flushleft}
\subsection*{Page 3}
\begin{flushleft}
\bigskip\textbf{Question 5.} A quadrilateral drawn in the complex plane has vertices A, B, C and D, labelled anticlock-\par
wise. These vertices are represented, respectively, by the complex numbers a, b, c and d.\par
Show that ABCD is a parallelogram (defined as a quadrilateral in which opposite sides are\par
parallel and equal in length) if and only if a + c = b + d . Show further that, in this case,\par
\[ABCD is a square if and only if i(a c) = b d.\]
- -\par
Let P QRS be a quadrilateral in the complex plane, with vertices labelled anticlockwise, the\par
\textasciicircum{}\circ\par
internal angles of which are all less than 180 . Squares with centres X, Y , Z and T are\par
constructed externally to the quadrilateral on the sides P Q, QR, RS and SP , respectively.\par
\textbf{(i)} If P and Q are represented by the complex numbers p and q, respectively, show that\par
X can be represented by\par
1 p(1 + i) + q(1 i) .\par
2 -\par
( )\par
\textbf{(ii)} Show that XYZT is a square if and only if P QRS is a parallelogram.\par
6 Starting from the result that\par
x\par
\[h(t) > 0 for 0 < t < x = h(t) dt > 0 ,\]
⇒ ∫\par
′′ ′\par
\[show that, if f (t) > 0 for 0 < t < x and f(0) = f (0) = 0, then f(t) > 0 for 0 < t < x .\]
0 0\par
\textbf{(i)} Show that, for 0 < x < 1 \pi,\par
\[cos x cosh x < 1 .\]
\textbf{(ii)} Show that, for 0 < x < 1 \pi,\par
1 sin x x\par
\[< < .\]
cosh x x sinh x\par
\end{flushleft}
\subsection*{Page 4}
\begin{flushleft}
\bigskip\textbf{Question 7.} The four distinct points P (i = 1, 2, 3, 4) are the vertices, labelled anticlockwise, of a cyclic\par
i\par
quadrilateral. The lines P P and P P intersect at Q.\par
1 3 2 4\par
\textbf{(i)} By considering the triangles P QP and P QP show that (P Q)(QP ) = (P Q)(QP ) .\par
1 4 2 3 1 3 2 4\par
\textbf{(ii)} Let p be the position vector of the point P (i = 1, 2, 3, 4). Show that there exist\par
i i\par
numbers a , not all zero, such that\par
i\par
4 4\par
\[a = 0 and a p = 0 . ( )\]
i i i\par
∗\par
\[∑i=1 ∑i=1\]
\textbf{(iii)} Let a (i = 1, 2, 3, 4) be any numbers, not all zero, that satisfy ( ). Show that\par
i\par
∗\par
\[a + a = 0 and that the lines P P and P P intersect at the point with position\]
1 3 1 3 2 4\par
̸\par
vector\par
a p + a p\par
1 1 3 3\par
.\par
a + a\par
1 3\par
\[Deduce that a a (P P )2 = a a (P P )2 .\]
1 3 1 3 2 4 2 4\par
\bigskip\textbf{Question 8.} The numbers f(r) satisfy f(r) > f(r + 1) for r = 1, 2, . . . . Show that, for any non-negative\par
integer n,\par
kn+1-1\par
kn(k 1) f(kn+1) ⩽ f(r) ⩽ kn(k 1) f(kn)\par
- -\par
\[r∑=kn\]
where k is an integer greater than 1.\par
\textbf{(i)} By taking f(r) = 1/r, show that\par
2N+1-1\par
N + 1 1\par
⩽ ⩽ N + 1 .\par
2 r\par
\[∑r=1\]
\infty\par
Deduce that the sum 1 does not converge.\par
r\par
\[r∑=1\]
\textbf{(ii)} By taking f(r) = 1/r3, show that\par
\infty\par
⩽ 1 1 .\par
\[r^3 3\]
\[∑r=1\]
\textbf{(iii)} Let S(n) be the set of positive integers less than n which do not have a 2 in their\par
decimal representation and let σ(n) be the sum of the reciprocals of the numbers in\par
\[S(n), so for example σ(5) = 1 + 1 + 1 . Show that S(1000) contains 93 1 distinct\]
3 4 -\par
numbers.\par
\[Show that σ(n) < 80 for all n.\]
\end{flushleft}
\subsection*{Page 5}
\begin{flushleft}
\bigskip\textbf{Section B: Mechanics}\par
\bigskip\textbf{Question 9.} A particle of mass m is projected with velocity u. It is acted upon by the force mg due to\par
gravity and by a resistive force mkv, where v is its velocity and k is a positive constant.\par
-\par
Given that, at time t after projection, its position r relative to the point of projection is\par
given by\par
kt 1 + e\par
-kt\par
1 e\par
-kt\par
\[r = - g + - u ,\]
\[k^2 k\]
find an expression for v in terms of k, t, g and u. Verify that the equation of motion and\par
the initial conditions are satisfied.\par
\textasciicircum{}\circ\par
\[Let u = u cos α i + u sin α j and g = g j, where 0 < α < 90 , and let T be the time after\]
-\par
\[projection at which r . j = 0. Show that\]
kT\par
\[uk sin α = 1 g .\]
( 1 e-kT - )\par
-\par
Let β be the acute angle between v and i at time T . Show that\par
(ekT 1)g\par
\[tan β = - tan α .\]
uk cos α -\par
\[Show further that tan β > tan α (you may assume that sinh kT > kT ) and deduce that β > α.\]
\bigskip\textbf{Question 10.} Two particles X and Y , of equal mass m, lie on a smooth horizontal table and are connected\par
by a light elastic spring of natural length a and modulus of elasticity λ. Two more springs,\par
identical to the first, connect X to a point P on the table and Y to a point Q on the table.\par
The distance between P and Q is 3a.\par
Initially, the particles are held so that XP = a, Y Q = 1 a , and P XY Q is a straight line.\par
The particles are then released.\par
At time t, the particle X is a distance a + x from P and the particle Y is a distance a + y\par
from Q. Show that\par
\[d2x λ\]
\[m = (2x + y)\]
\[dt2 - a\]
\[d2y\]
and find a similar expression involving . Deduce that\par
\[dt2\]
\[x y = A cos ωt + B sin ωt\]
-\par
\[where A and B are constants to be determined and maω2 = λ. Find a similar expression for\]
x + y.\par
Show that Y will never return to its initial position.\par
\end{flushleft}
\subsection*{Page 6}
\begin{flushleft}
\bigskip\textbf{Question 11.} A particle P of mass m is connected by two light inextensible strings to two fixed points A\par
and B, with A vertically above B. The string AP has length x. The particle is rotating\par
about the vertical through A and B with angular velocity ω, and both strings are taut.\par
Angles P AB and P BA are α and β, respectively.\par
Find the tensions T and T in the strings AP and BP (respectively), and hence show that\par
A B\par
ω2x cos α ⩾ g.\par
\[Consider now the case that ω2x cos α = g. Given that AB = h and BP = d, where h > d,\]
\[show that h cos α ⩾ \sqrt\{h^2\} d^2. Show further that\]
-\par
mgh\par
\[mg < T ⩽ .\]
A\par
\[\sqrt\{h^2\} d^2\]
-\par
Describe the geometry of the strings when T attains its upper bound.\par
A\par
\end{flushleft}
\subsection*{Page 7}
\begin{flushleft}
\bigskip\textbf{Section C: Probability and Statistics}\par
\bigskip\textbf{Question 12.} The random variable X has probability density function f(x) (which you may assume is\par
\[differentiable) and cumulative distribution function F(x) where < x < . The random\]
-\infty \infty\par
\[variable Y is defined by Y = eX . You may assume throughout this question that X and Y\]
have unique modes.\par
\textbf{(i)} Find the median value y of Y in terms of the median value x of X.\par
m m\par
\textbf{(ii)} Show that the probability density function of Y is f(ln y)/y, and deduce that the\par
′\par
\[mode λ of Y satisfies f (ln λ) = f(ln λ).\]
\textbf{(iii)} Suppose now that X N(µ, σ2), so that\par
∼\par
\[f(x) = e\]
-(x-µ)2/(2σ2)\par
.\par
\[σ\sqrt\{2\}\pi\]
Explain why\par
\infty\par
e\par
-(x-µ-σ2)2/(2σ2)\par
\[dx = 1\]
\[σ\sqrt\{2\}\pi ∫-\infty\]
\[and hence show that E(Y ) = eµ+ 2 1 σ2 .\]
\textbf{(iv)} Show that, when X N(µ, σ2),\par
∼\par
\[λ < y < E(Y ) .\]
m\par
\end{flushleft}
\subsection*{Page 8}
\begin{flushleft}
\bigskip\textbf{Question 13.} I play a game which has repeated rounds. Before the first round, my score is 0. Each round\par
can have three outcomes:\par
1. my score is unchanged and the game ends;\par
2. my score is unchanged and I continue to the next round;\par
3. my score is increased by one and I continue to the next round.\par
The probabilities of these outcomes are a, b and c, respectively (the same in each round),\par
\[where a + b + c = 1 and abc = 0. The random variable N represents my score at the end of\]
̸\par
a randomly chosen game.\par
Let G(t) be the probability generating function of N .\par
\textbf{(i)} Suppose in the first round, the game ends. Show that the probability generating\par
function conditional on this happening is 1.\par
\textbf{(ii)} Suppose in the first round, the game continues to the next round with no change in\par
score. Show that the probability generating function conditional on this happening\par
is G(t).\par
\textbf{(iii)} By comparing the coefficients of tn, show that G(t) = a + bG(t) + ctG(t) . Deduce that,\par
for n ⩾ 0,\par
acn\par
\[P (N = n) = .\]
(1 b)n+1\par
-\par
\textbf{(iv)} Show further that, for n ⩾ 0,\par
µn\par
\[P (N = n) = ,\]
(1 + µ)n+1\par
\[where µ = E(N ).\]
\end{flushleft}
\end{document}
